\section{Hệ thống dữ liệu}\label{sec:du-lieu-nguon}
Dữ liệu của dự án đến từ hai bộ dữ liệu công khai của Olist trên Kaggle, được mô phỏng thành hai hệ thống nguồn độc lập trên PostgreSQL. Hệ thống thứ nhất là hệ thống bán hàng (OLTP), gồm 9 bảng với 1.550.922 dòng, ghi nhận đơn hàng, sản phẩm đã bán, thanh toán, đánh giá, khách hàng, người bán, danh mục sản phẩm và toạ độ địa lý theo mã bưu chính. Hệ thống thứ hai là hệ thống CRM/Marketing, gồm 2 bảng với 8.842 dòng: 8.000 khách tiềm năng (MQL) và 842 thương vụ đã chốt. Hai hệ thống chỉ gặp nhau qua mã người bán \code{seller\_\allowbreak{}id}; vì là hai hệ thống khác nhau nên cơ sở dữ liệu không đặt khoá ngoại giữa chúng và việc liên kết được thực hiện ở kho dữ liệu thông qua chiều \code{dim\_\allowbreak{}seller}. Tổng dung lượng 11 file CSV là 121,18 MB (1.559.764 dòng), riêng bảng toạ độ địa lý chiếm 58,44 MB và hơn một triệu dòng. Mỗi hệ thống nằm trong một schema riêng (\code{oltp}, \code{crm}), giữ nguyên tên cột gốc, kể cả hai cột sai chính tả \code{product\_\allowbreak{}name\_\allowbreak{}lenght} và \code{product\_\allowbreak{}description\_\allowbreak{}lenght}, vì đó là đặc điểm của hệ thống nguồn và chỉ được sửa ở tầng ODS. Dữ liệu được nạp nguyên trạng, không làm sạch. Nhóm tạo hai database: \code{olist\_\allowbreak{}full} (262,01 MB) chứa toàn bộ dữ liệu để khám phá, và \code{olist\_\allowbreak{}dw} (185,14 MB) chỉ chứa các bản ghi đến hết 31/12/2017 (1.273.620 dòng) để làm điểm xuất phát cho nạp tăng dần. Số dòng nạp vào \code{olist\_\allowbreak{}full} khớp hoàn toàn với số dòng trong file CSV và không có dòng nào bị loại.

\begin{table}[htbp]
\centering\small
\caption{Hai hệ thống nguồn: số bảng ghi, số cột và dung lượng file CSV}
\label{tab:src-tables}
\begin{tabular}{llrrr}
\toprule
\textbf{Bảng nguồn} & \textbf{File CSV} & \textbf{Số dòng} & \textbf{Số cột} & \textbf{MB} \\
\midrule
oltp.product\_category\_name\_translation & product\_category\_name\_translation.csv & 71 & 2 & 0,00 \\
oltp.products & olist\_products\_dataset.csv & 32.951 & 9 & 2,27 \\
oltp.sellers & olist\_sellers\_dataset.csv & 3.095 & 4 & 0,17 \\
oltp.geolocation & olist\_geolocation\_dataset.csv & 1.000.163 & 6 & 58,44 \\
oltp.customers & olist\_customers\_dataset.csv & 99.441 & 5 & 8,62 \\
oltp.orders & olist\_orders\_dataset.csv & 99.441 & 8 & 16,84 \\
oltp.order\_items & olist\_order\_items\_dataset.csv & 112.650 & 7 & 14,72 \\
oltp.order\_payments & olist\_order\_payments\_dataset.csv & 103.886 & 5 & 5,51 \\
oltp.order\_reviews & olist\_order\_reviews\_dataset.csv & 99.224 & 7 & 13,78 \\
crm.mql & olist\_marketing\_qualified\_leads\_dataset.csv & 8.000 & 4 & 0,67 \\
crm.closed\_deals & olist\_closed\_deals\_dataset.csv & 842 & 14 & 0,16 \\
Tổng (11 bảng) &  & 1.559.764 &  & 121,18 \\
\bottomrule
\end{tabular}
\end{table}

Hai hệ thống nguồn được mô phỏng thành hai schema riêng (\code{oltp} và \code{crm}) trong PostgreSQL; sau khi nạp đủ, kho có số dòng từng lớp như Bảng \ref{tab:layers}.
\begin{table}[htbp]
\centering\small
\caption{Số bảng và số dòng của từng lớp trong kho dữ liệu (trạng thái cuối)}
\label{tab:layers}
\begin{tabular}{llrr}
\toprule
\textbf{Lớp} & \textbf{Schema} & \textbf{Số bảng} & \textbf{Số dòng} \\
\midrule
Nguồn bán hàng & oltp & 9 & 1.550.922 \\
Nguồn marketing & crm & 2 & 8.842 \\
Staging & staging & 11 & 1.559.764 \\
ODS & ods & 18 & 678.183 \\
Kho (dimension + fact) & dw & 16 & 614.223 \\
Giám sát ETL & etl\_log & 6 & 2.701 \\
Data mart (view và materialized view) & mart\_* & 17 view & -- \\
\bottomrule
\end{tabular}
\end{table}

\paragraph{Tần suất và khối lượng cập nhật.} Dữ liệu Kaggle là ảnh chụp tĩnh nên nhóm mô phỏng cập nhật theo tháng: nạp ban đầu đến 31/12/2017 (45.430 đơn) rồi 11 đợt hằng tháng (01--11/2018), tổng 286.144 bản ghi mới (xem mục 
ef{sec:doi-soat}).
\todo{Nhóm nêu rõ tần suất cập nhật giả định của hệ thống thực (ví dụ đơn hàng cập nhật liên tục, kho cập nhật hằng ngày) và dung lượng mỗi lần cập nhật.}

\subsection{Tần suất cập nhật và công nghệ xử lý}\label{sec:tan-suat}
Hình \ref{fig:tan-suat} cho thấy tần suất dữ liệu mới theo thời gian: số đơn mỗi ngày, số đơn và số dòng hàng mỗi tháng, số review, số khách tiềm năng (MQL) và số deal chốt mỗi tháng. Trong giai đoạn 01/2017--08/2018 (608 ngày, chứa 99,6\% số đơn), hệ thống nhận trung bình 163,0 đơn mỗi ngày với 184,7 dòng hàng và 170,3 lần thanh toán; ngày cao nhất là Black Friday 24/11/2017 với 1.176 đơn, gấp 7,2 lần mức trung bình. Phía marketing nhỏ hơn nhiều: trung bình 13,2 MQL và 1,3 deal mỗi ngày. Số review theo ngày tạo review bám theo số đơn với độ trễ vài ngày đến hai tuần, nên đỉnh review (463 review ngày 19/12/2017) đến sau đỉnh đơn.
\hinh{2-4_tan_suat_cap_nhat.png}{Tần suất cập nhật của các bảng nguồn theo thời gian}{fig:tan-suat}{1.0}
Với khối lượng này, bảng \ref{tab:tan-suat} đề xuất công nghệ xử lý cho từng bảng. Đơn hàng đổi trạng thái liên tục nhưng chỉ vài trăm đơn mỗi ngày, nên lô hằng ngày bằng procedure PostgreSQL (đang dùng trong dự án; mỗi đợt khoảng 20 giây) là đủ. Khi cần dashboard gần thời gian thực thì thêm CDC kết hợp Kafka giữa hệ thống nguồn và staging (nối sang chương 5). Hai bảng tham chiếu (toạ độ, bản dịch danh mục) gần như tĩnh nên chỉ nạp lại theo quý.
\begin{table}[htbp]\centering\scriptsize\caption{Tần suất cập nhật và công nghệ xử lý đề xuất cho từng bảng nguồn (cửa sổ 01/2017--08/2018)}\label{tab:tan-suat}
\begin{tabular}{p{2.66cm}p{1.38cm}p{2.02cm}p{3.31cm}p{4.23cm}}\toprule
\textbf{Bảng nguồn} & \textbf{Trung bình / ngày} & \textbf{Đỉnh một ngày} & \textbf{Tần suất cập nhật} & \textbf{Công nghệ xử lý đề xuất}\\\midrule
\code{oltp.orders} & 163,0 & 1.176 (24/11/2017) & liên tục: mỗi đơn đổi trạng thái 4–5 lần (mua, duyệt, bàn giao, giao) & lô hằng ngày bằng procedure PostgreSQL; nếu cần gần thời gian thực: CDC + Kafka (chương 5)\\
\code{oltp.order\_items} & 184,7 & 1.366 (24/11/2017) & cùng lúc với đơn (ghi một lần khi đặt) & lô hằng ngày, nạp theo khoá (order\_id, order\_item\_id)\\
\code{oltp.order\_payments} & 170,3 & 1.214 (24/11/2017) & cùng lúc với đơn & lô hằng ngày; có thể gộp chung job với đơn\\
\code{oltp.order\_reviews} & 162,7 & 463 (19/12/2017) & sau khi giao hàng (trễ vài ngày đến hai tuần) & lô hằng ngày; thử lại dòng mồ côi khi đơn chưa về\\
\code{oltp.customers} & 157,5 & 1.143 (24/11/2017) & khi khách mới đặt đơn đầu tiên & lô hằng ngày; SCD2 khi đổi mã bưu chính\\
\code{oltp.sellers} & 4,9 & 24 (26/01/2017) & khi seller bán đơn đầu tiên & lô hằng ngày\\
\code{oltp.products} & 53,7 & 374 (24/11/2017) & khi sản phẩm bán lần đầu & lô hằng ngày\\
\code{crm.mql} & 13,2 & 93 (02/05/2018) & hằng ngày (khách tiềm năng để lại thông tin) & lô hằng ngày; luồng sự kiện nếu cần phản hồi nhanh cho đội sales\\
\code{crm.closed\_deals} & 1,3 & 22 (30/04/2018) & vài ngày một lần (khi chốt hợp đồng) & lô hằng ngày\\
\code{oltp.geolocation} & 0 & 0 & tĩnh (bảng tham chiếu, đổi khi có bản đồ mới) & nạp lại toàn phần theo quý; gom về một điểm cho mỗi mã bưu chính\\
\code{oltp.product\_category\_name\_translation} & 0 & 0 & tĩnh (71 dòng) & bảng seed, nạp lại khi có danh mục mới\\
\bottomrule\end{tabular}\end{table}


\subsection{Phân loại dữ liệu (data taxonomy)}\label{sec:taxonomy}
Mười một bảng nguồn được chia thành bốn nhóm (Hình \ref{fig:taxonomy}): \textbf{dữ liệu chủ} (khách, seller, sản phẩm: mô tả đối tượng, ít đổi), \textbf{dữ liệu giao dịch} (đơn, dòng hàng, thanh toán: ghi sự kiện mua bán, tăng liên tục), \textbf{dữ liệu tương tác} (review, MQL, deal: phản hồi của khách và hành trình thu hút seller) và \textbf{dữ liệu tham chiếu} (toạ độ, bản dịch danh mục: gần như tĩnh). Mỗi bảng kèm giá trị thật của hai cột tiêu biểu. Bảng \ref{tab:taxonomy} phân loại thêm theo kiểu dữ liệu, độ nhạy cảm (định danh, vị trí) và tính biến động; kết quả dùng để quyết định tần suất nạp và cách xử lý: bảng có tính biến động cao (đơn hàng) cần nạp tăng dần theo watermark, bảng tĩnh chỉ cần nạp lại toàn phần.
\hinh{2-4_data_taxonomy.png}{Cây phân loại dữ liệu nguồn (data taxonomy) với ví dụ giá trị thật}{fig:taxonomy}{1.0}
\begin{table}[htbp]\centering\scriptsize\caption{Phân loại 11 bảng nguồn theo kiểu dữ liệu, độ nhạy cảm và tính biến động}\label{tab:taxonomy}
\begin{tabular}{p{2.57cm}p{2.11cm}p{3.03cm}p{3.58cm}p{2.30cm}}\toprule
\textbf{Bảng} & \textbf{Nhóm} & \textbf{Kiểu dữ liệu} & \textbf{Độ nhạy cảm} & \textbf{Tính biến động}\\\midrule
\code{customers} & Dữ liệu chủ (master) & định danh (khách, vị trí thành phố/bang) & nhạy cảm vừa: mã khách đã ẩn danh, có vị trí & thấp (đổi khi khách đổi địa chỉ)\\
\code{sellers} & Dữ liệu chủ (master) & định danh (seller, vị trí) & nhạy cảm vừa: mã seller ẩn danh, có vị trí & thấp\\
\code{products} & Dữ liệu chủ (master) & thuộc tính (danh mục, kích thước) & không nhạy cảm & thấp\\
\code{orders} & Dữ liệu giao dịch & thời gian và trạng thái & nhạy cảm thấp: gắn với khách qua customer\_id & cao (trạng thái đổi theo từng mốc)\\
\code{order\_items} & Dữ liệu giao dịch & số tiền và khoá & thông tin thương mại (giá, phí) & thấp (ghi một lần)\\
\code{order\_payments} & Dữ liệu giao dịch & số tiền, hình thức & thông tin thanh toán (đã tổng hợp, không có số thẻ) & thấp\\
\code{order\_reviews} & Dữ liệu tương tác & điểm số và văn bản tự do & nhạy cảm vừa: văn bản có thể nêu thông tin cá nhân & thấp (hiếm khi sửa)\\
\code{mql} & Dữ liệu tương tác & ngày, kênh (phân loại) & thông tin kinh doanh nội bộ & thấp\\
\code{closed\_deals} & Dữ liệu tương tác & phân loại và số liệu doanh nghiệp seller & thông tin kinh doanh nội bộ & thấp\\
\code{geolocation} & Dữ liệu tham chiếu & vị trí (vĩ độ, kinh độ) & vị trí: chỉ ở mức mã bưu chính & rất thấp (tĩnh)\\
\code{category\_translation} & Dữ liệu tham chiếu & văn bản tra cứu & không nhạy cảm & rất thấp (tĩnh)\\
\bottomrule\end{tabular}\end{table}


\subsection{Luồng dữ liệu giữa các hệ thống CNTT}\label{sec:dataflow}
Hình \ref{fig:flow-now} mô tả hiện trạng: mũi tên ghi \emph{thông tin được truyền} (thông tin đơn, trạng thái giao, thông tin thanh toán, điểm đánh giá, thông tin lead), không ghi hành động. Hệ thống bán hàng và hệ thống CRM là hai hệ thống độc lập, chỉ gặp nhau qua mã người bán, nên mỗi báo cáo hiện lấy số trực tiếp từ từng hệ thống và khó đối chiếu với nhau.
\hinh{2-4_data_flow_hien_trang.png}{Luồng dữ liệu giữa các hệ thống CNTT hiện trạng (bản nháp)}{fig:flow-now}{1.0}
Hình \ref{fig:flow-new} là luồng đề xuất: nguồn $\rightarrow$ staging $\rightarrow$ ODS $\rightarrow$ kho dữ liệu $\rightarrow$ data mart $\rightarrow$ BI, kèm lớp giám sát \code{etl\_log}; Bảng \ref{tab:p8-dataflow} nêu thành phần, công nghệ và tần suất của từng lớp.
\hinh{2-4_data_flow_de_xuat.png}{Luồng dữ liệu đề xuất từ nguồn đến BI (bản nháp)}{fig:flow-new}{1.0}
\begin{table}[htbp]\centering\scriptsize\caption{Luồng dữ liệu đề xuất: lớp, thành phần, công nghệ, tần suất và ghi chú}\label{tab:p8-dataflow}
\begin{tabular}{p{1.43cm}p{3.22cm}p{3.22cm}p{2.86cm}p{2.86cm}}\toprule
\textbf{Lớp} & \textbf{Thành phần} & \textbf{Công nghệ} & \textbf{Tần suất} & \textbf{Ghi chú}\\\midrule
Nguồn & oltp (9 bảng), crm (2 bảng) & PostgreSQL 18; nạp từ 11 file CSV bằng Python (COPY) & mô phỏng cập nhật theo tháng (ban đầu đến 31/12/2017, rồi 11 đợt) & hai hệ thống độc lập, chỉ gặp nhau qua seller\_id\\
Staging & stg\_* (bản sao thô) & PostgreSQL; extract bằng Python + SQL & mỗi lần chạy pipeline (full hoặc incremental) & thêm mã lô nạp và thời điểm nạp để truy vết\\
ODS & ods.* (dữ liệu sạch, cờ chất lượng) & PostgreSQL procedure (PL/pgSQL) & mỗi lần chạy pipeline & dòng lỗi ghi \code{rejected\_rows}, không xoá âm thầm\\
Kho dữ liệu & dw: 10 dimension + 6 fact & PostgreSQL procedure; khoá thay thế, SCD1/SCD2 & mỗi lần chạy pipeline; nạp tăng dần theo watermark & 11 view đóng vai cho \code{dim\_date}\\
Data mart & mart\_sales, mart\_fulfillment, mart\_cx, mart\_acquisition & PostgreSQL view và materialized view & làm mới materialized view mỗi lần chạy & tổng hợp sẵn cho dashboard DB1--DB4\\
Giám sát & etl\_log: job\_run, rejected\_rows, dq\_check... & PostgreSQL; Python ghi log & ghi ở mọi bước & nguồn cho bảng giám sát ở mục 3.3\\
BI & 37 file CSV, mô hình ngữ nghĩa, 4 dashboard & Python xuất CSV; Power BI Desktop (Import, DAX) & làm mới sau mỗi lần xuất CSV & Publish to web\\
\bottomrule\end{tabular}\end{table}

\todo{Nhóm vẽ lại hai sơ đồ luồng dữ liệu bằng draw.io từ docs/diagrams/data-flow-hien-trang.drawio và data-flow-de-xuat.drawio (giữ tên file hình); kiểm tra tên các hệ thống ngoài (marketplace, vận chuyển, thanh toán) cho đúng cách hiểu nghiệp vụ của nhóm.}
